\subsection{An explicit eight-torsion ambiguity in the unramified dyadic case}

\begin{proposition}[B-RD43: even eight-torsion does not suffice at $p=2,e=1$]
\label{prop:B-RD43}
\status{Proved. Complete root and independent original-source, proof and exact-certificate audits passed.}
Let $K=\mathbf Q_2(\omega)$, where $\omega^2+\omega+1=0$, and put
$\lambda=1+2\omega$, so $\lambda^2=-3$. Define
\[
 E:\ y^2+(5-\lambda)xy-4\lambda y=x^3-2\lambda x^2,
 \qquad P=(0,0),\qquad B=E/\langle P\rangle,
 \qquad q=-\frac{4\lambda}{3}.
                                                        \tag{B-RD43.1}
\]
Then $P$ has exact order eight, $v_K(q)=2>0$, and the split Tate
curve $T_q$ satisfies
\[
                             B[8]\simeq T_q[8].
                                                        \tag{B-RD43.2}
\]
The requested characters of $B[2^\infty]$ and $T_q[2^\infty]$
differ on an inertia element whose image has order three in the
finite special-inertia quotient. Consequently
\[
                       \nu_{\min}(2,1,h)\ge3\qquad(h\ge2).
                                                        \tag{B-RD43.3}
\]
The same lower bound holds for every absolute ramification index
$e$ not divisible by three. No assertion of sufficiency of
sixteen-torsion is made.
\end{proposition}
\begin{proof}
The field $K$ is unramified quadratic over $\mathbf Q_2$, since
$X^2+X+1$ is separable and irreducible over $\mathbf F_2$.
The displayed equation is integral. Exact addition gives
\[
\begin{array}{c|c|c}
 n&x(nP)&y(nP)\\\hline
 1&0&0\\
 2&2\lambda&-6-6\lambda\\
 3&-6+2\lambda&12\\
 4&(-3+\lambda)/2&3\\
 5&-6+2\lambda&12-12\lambda\\
 6&2\lambda&0\\
 7&0&4\lambda
\end{array}
\qquad 8P=0.
                                                        \tag{B-RD43.4}
\]
In particular $P$ has exact order eight. This curve arises from
the $X_1(8)$ Tate normal form in
\citet[Appendix~A.2, pp.~26--27, Figure~8]{AntieauAuel2021Descent}
by taking its parameter to be $\lambda$ and scaling
$x=X/4$, $y=Y/8$. The displayed addition certificate verifies
its torsion without relying on that parametrization.

We compute the exact extension class in
Lemma~\ref{lem:B-RD29-isogeny-Tate}, so that the Tate parameter
in (B-RD43.1) is specified rather than merely asserted to exist.
The tangent at $P$ is $y=0$. The tangent at $2P$ is
\[
 \ell_2=y-\frac{-9+\lambda}{2}x-(-3+3\lambda)=0.
\]
Put $x_2=2\lambda$ and $x_4=(-3+\lambda)/2$.
The Miller doubling identity gives a function with divisor
$8[P]-8[0]$:
\[
 f=\frac{y^4\ell_2^2}{(x-x_2)^4(x-x_4)}.
\]
Here the last tangent is the vertical line at the two-torsion point
$4P$. Reducing by the equation of $E$ cancels the denominator and gives
\[
\begin{split}
 f(x,y)={}&x^4-4\lambda x^3-8x^2
            +(2\lambda-6)x^2y\\
          &+(12\lambda+4)xy+16y.
\end{split}                                             \tag{B-RD43.5}
\]
This polynomial identity and all the displayed point additions
are verified exactly in $\mathbf Q(\lambda)$.
The divisor description of the dual-isogeny connecting homomorphism
is \citet[Appendix~A.1, Lemmas~A.1--A.2, p.~22]
{AntieauAuel2021Descent}. Its unspecified normalizing constant
cancels in the ratio at $3P$ and $2P$, since both avoid the zero
and pole of $f$. Thus
\[
 \delta(P)=\frac{f(3P)}{f(2P)}
           =\frac{576\lambda}{-432}
           =-\frac{4\lambda}{3}
                   \quad\hbox{in }K^\times/(K^\times)^8.
                                                        \tag{B-RD43.6}
\]
The quotient $B=E/\langle P\rangle$ therefore fits into the
Kummer extension
$0\to\mu_8\to B[8]\to\mathbf Z/8\to0$ with the indicated
class. Since $v_K(q)=2$, Lemma~\ref{lem:B-RD29-isogeny-Tate}
and Tate's explicit finite-order construction identify it with
$T_q[8]$. As an additional check, the Figure~8 representative
$\lambda^7(\lambda-1)^6(\lambda+1)^4$ divided by (B-RD43.6)
is exactly $(-3-\lambda)^8$.

The invariant calculation gives
\[
 c_4=-200-312\lambda,\qquad
 c_6=15488+576\lambda,\qquad
 \Delta=20736(\lambda-2),
\]
so
\[
                     v_K(c_4)=4,\quad v_K(\Delta)=8,
                     \quad v_K(j)=4.
                                                        \tag{B-RD43.7}
\]
For example, $c_4=-16(32+39\omega)$ and
$\lambda-2=-1+2\omega$ are respectively sixteen times a unit
and a unit. The displayed discriminant is minimal, since eight
is less than twelve and integral discriminant valuations are
nonnegative. Integral $j$ gives potential good reduction by
\citet[Section~2, p.~496]{SerreTate1968Good}; its positive valuation
gives supersingular special reduction in characteristic two,
as in the already verified proof of Proposition~\ref{prop:B-RD22}.

Over the completed maximal unramified field, let $m$ be a minimal
good-reduction degree. The change to a unit-discriminant equation
implies $12\mid8m$, hence $3\mid m$. Faithful special inertia
therefore contains an element of order three. Its trace on the
rank-two rational Dieudonne module is $-1$, as in
Proposition~\ref{prop:B-RD22}. The isogeny $E\to B$ preserves this
rational module and its finite inertia action. The Tate group's
two good rank-one graded pieces have trace $2$ on that element.
The order-three assertion concerns the finite special-inertia image;
no finite-order lift in the absolute inertia group is asserted.
Both groups satisfy the manuscript monodromy condition, and the
characters differ despite (B-RD43.2).

Finally, take an Eisenstein extension of $K$ of degree $e$.
The torsion isomorphism persists and the raw discriminant valuation
becomes $8e$. If $3\nmid e$, the relation $12\mid8em$ still
forces $3\mid m$, so the same character discrepancy survives.
Adding $h-2$ identical constant etale factors proves every height
$h\ge2$. Since eight-torsion is $M(2)$, its failure proves
(B-RD43.3).
\end{proof}

\paragraph{Exact certificate and scope.}
The program \path{code/verify_p2_character_level3.py} checks all eight
subgroup points and all $64$ additions, the exact invariant
polynomials, both tangent lines, the affine Miller-function identity,
its two evaluations and the eighth-power parameter comparison.
Its retained author output records $0.5159330129972659$ seconds.
The finite-scheme extension dictionary and the special-inertia
argument are proved above using the stated primary sources.
This counterexample raises the lower bound; it does not compute
the least dyadic index.

Independent author-output-preserving reruns of the full certificate passed
in $0.5299817019986222$ and $0.6142895730008604$ seconds. Both reviewers
checked the complete proof and the original descent and Tate sources.
